% c14 原书 p.168–180（PDF p.175–187）
\begingroup
\setlength{\emergencystretch}{2.5em}
\subsection{\v{C}ech 上同调}
\label{sec:2.3}

设 $X$ 是拓扑空间，$\fU=\{U_i\}_{i\in I}$ 是 $X$ 的开覆盖。对任意 $i_0,\dots,i_n\in I$，令 $U_{i_0\dots i_n}=U_{i_0}\cap\cdots\cap U_{i_n}$。给定 $X$ 上交换群的前层 $\cP$，如下定义复形 $C^\cdot(\fU,\cP)$：对每个 $n\ge 0$，令
\[
C^n(\fU,\cP)=\prod_{i_0,\dots,i_n\in I}\cP(U_{i_0\dots i_n}),
\]
并定义
\[
d\colon C^n(\fU,\cP)\to C^{n+1}(\fU,\cP),\quad s=(s_{i_0\dots i_n})\mapsto ds=((ds)_{i_0\dots i_{n+1}})
\]
为
\[
(ds)_{i_0\dots i_{n+1}}=\sum_{k=0}^{n+1}(-1)^k s_{i_0\dots\hat{i_k}\dots i_{n+1}}|_{U_{i_0\dots i_{n+1}}},
\]
其中 $s_{i_0\dots\hat{i_k}\dots i_{n+1}}|_{U_{i_0\dots i_{n+1}}}$ 是 $s_{i_0\dots\hat{i_k}\dots i_{n+1}}\in \cP(U_{i_0\dots\hat{i_k}\dots i_{n+1}})$ 在限制映射
\[
\cP(U_{i_0\dots\hat{i_k}\dots i_{n+1}})\to \cP(U_{i_0\dots i_{n+1}})
\]
下的像。可以验证 $dd=0$。我们称 $C^\cdot(\fU,\cP)$ 为 $\cP$ 关于开覆盖 $\fU$ 的\emph{\v{C}ech 复形}\idx{cechfuxing}{\v{C}ech 复形}。$C^\cdot(\fU,\cP)$ 的第 $n$ 个上同调群记作 $\check{H}^n(\fU,\cP)$，称为 $\cP$ 关于开覆盖 $\fU$ 的第 $n$ 个\emph{\v{C}ech 上同调群}\idx{cechshangtongdiaoqun}{\v{C}ech 上同调群}。

固定 $I$ 上的一个全序。如下定义复形 $C'^\cdot(\fU,\cP)$：对每个 $n$，令
\[
C'^n(\fU,\cP)=\prod_{i_0<\cdots<i_n}\cP(U_{i_0\dots i_n}),
\]
并用与上面相同的公式定义 $d\colon C'^n(\fU,\cP)\to C'^{n+1}(\fU,\cP)$。我们称 $C'^\cdot(\fU,\cP)$ 为 $\cP$ 关于开覆盖 $\fU$ 的\emph{交错 \v{C}ech 复形}\idx{jiaocuocechfuxing}{交错 \v{C}ech 复形}。

\begin{proposition}\label{it:2.3.1}
设 $X$ 是拓扑空间，$\fU=\{U_i\}_{i\in I}$ 是 $X$ 的开覆盖，$\cP$ 是 $X$ 上交换群的前层。则存在复形的态射 $f^\cdot\colon C'^\cdot(\fU,\cP)\to C^\cdot(\fU,\cP)$，它在上同调群上诱导出同构
\[
H^n(C'^\cdot(\fU,\cP))\cong H^n(C^\cdot(\fU,\cP)).
\]
\end{proposition}

\begin{proof}
对每个 $n\ge 0$，设 $K_n(I)$ 是以 $(i_0,\dots,i_n)\in I^{n+1}$ 为基的自由交换群，$K'_n(I)$ 是 $K_n(I)$ 中由满足 $i_0<\cdots<i_n$ 的 $(i_0,\dots,i_n)$ 生成的自由交换子群。定义
\[
\partial_n\colon K_n(I)\to K_{n-1}(I)
\]
为
\[
\partial_n(i_0,\dots,i_n)=\sum_{k=0}^n(-1)^k(i_0,\dots,\hat{i_k},\dots,i_n).
\]

则 $\partial_{n-1}\partial_n=0$，且 $\partial_n(K'_n(I))\subset K'_{n-1}(I)$ 对任意 $n$ 成立。于是 $K_\cdot(I)$ 是复形，$K'_\cdot(I)$ 是子复形。（同态 $\partial$ 的次数为 $-1$。）设 $i_\cdot\colon K'_\cdot(I)\to K_\cdot(I)$ 是包含映射。如下定义同态 $f_n\colon K_n(I)\to K'_n(I)$（$n\ge 0$）：
\[
f_n(i_0,\dots,i_n)=
\begin{cases}
0, & \text{若存在 } k\ne l \text{ 使 } i_k=i_l,\\
\operatorname{sgn}(\sigma)(i_{\sigma(0)},\dots,i_{\sigma(n)}), & \text{否则},
\end{cases}
\]
其中 $\sigma$ 是 $\{0,\dots,n\}$ 的置换，满足 $i_{\sigma(0)}<\cdots<i_{\sigma(n)}$，而 $\operatorname{sgn}(\sigma)$ 是 $\sigma$ 的符号。我们留给读者证明 $f_\cdot\colon K_\cdot(I)\to K'_\cdot(I)$ 是复形的态射。我们有
\[
f_\cdot i_\cdot=\id_{K'_\cdot(I)}.
\]
我们断言，存在态射 $h_n\colon K_n(I)\to K_{n+1}(I)$（$n\ge 0$），使得
\[
\partial_{n+1}h_n+h_{n-1}\partial_n=i_nf_n-\id_{K_n(I)},
\]
并且对任意 $(i_0,\dots,i_n)\in I^{n+1}$，$h_n(i_0,\dots,i_n)$ 是那些满足 $\{j_0,\dots,j_{n+1}\}\subset\{i_0,\dots,i_n\}$（作为集合）的 $(j_0,\dots,j_{n+1})\in I^{n+2}$ 的线性组合。先承认这一断言。设
\[
h_n(i_0,\dots,i_n)=\sum_{\{j_0,\dots,j_{n+1}\}\subset\{i_0,\dots,i_n\}}c^{i_0\dots i_n}_{j_0\dots j_{n+1}}(j_0,\dots,j_{n+1}).
\]
我们定义
\[
h^*_n\colon C^{n+1}(\fU,\cP)\to C^n(\fU,\cP),\quad s=(s_{j_0\dots j_{n+1}})\mapsto h^*_n(s)=((h^*_n(s))_{i_0\dots i_n})
\]
为
\[
(h^*_n(s))_{i_0\dots i_n}=\sum_{\{j_0,\dots,j_{n+1}\}\subset\{i_0,\dots,i_n\}}c^{i_0\dots i_n}_{j_0\dots j_{n+1}}s_{j_0\dots j_{n+1}}|_{U_{i_0\dots i_n}}.
\]
类似地定义 $i^*_n\colon C^n(\fU,\cP)\to C'^n(\fU,\cP)$、$f^*_n\colon C'^n(\fU,\cP)\to C^n(\fU,\cP)$ 与 $\partial^*_{n+1}\colon C^n(\fU,\cP)\to C^n(\fU,\cP)$。可以证明 $i^*$ 与 $f^*$ 是复形的态射，并且
\begin{align*}
\partial^*_{n+1}&=d_n,\\
i^*_\cdot f^*_\cdot&=(f_\cdot i_\cdot)^*=\id_{C'^\cdot(\fU,\cP)},\\
d_{n-1}h^*_{n-1}+h^*_nd_n&=(\partial_{n+1}h_n+h_{n-1}\partial_n)^*\\
&=(i_nf_n-\id_{K_n(I)})^*\\
&=f^*_ni^*_n-\id_{C^n(\fU,\cP)}.
\end{align*}
这说明 $f^*$ 与 $i^*$ 在上同调群上诱导的同态互为逆映射。

在证明我们的断言之前，先证明：对任意集合 $I$ 和任意 $n\ge 1$，有 $H_n(K_\cdot(I))=0$。

固定 $I$ 中的一个元素 $i$。考虑同态
\[
g_n\colon K_n(I)\to K_{n+1}(I),\quad (i_0,\dots,i_n)\mapsto (i,i_0,\dots,i_n).
\]
可以验证
\[
\partial_{n+1}g_n+g_{n-1}\partial_n=\id_{K_n(I)}
\]
对所有 $n\ge 1$ 成立。由此得到 $H_n(K_\cdot(I))=0$ 对所有 $n\ge 1$ 成立。

下面证明我们的断言。注意 $i_0f_0\colon K_0(I)\to K_0(I)$ 与 $\id_{K_0(I)}$ 相同。我们定义 $h_0\colon K_0(I)\to K_1(I)$ 为零。设已对任意 $0\le k<n$ 定义了 $h_k\colon K_k(I)\to K_{k+1}(I)$，使得
\[
\partial_{k+1}h_k+h_{k-1}\partial_k=i_kf_k-\id_{K_k(I)},
\]
并且对任意 $(i_0,\dots,i_k)\in I^{k+1}$，$h_k(i_0,\dots,i_k)$ 是那些满足\allowbreak $\{j_0,\dots,j_{k+1}\}\subset\{i_0,\dots,i_k\}$ 的 $(j_0,\dots,j_{k+1})\in I^{k+2}$ 的线性组合。于是对任意 $(i_0,\dots,i_n)\in I^{n+1}$，\allowbreak $(i_nf_n-\id_{K_n(I)}-h_{n-1}\partial_n)$\allowbreak $(i_0,\dots,i_n)$ 是那些满足\allowbreak $\{j_0,\dots,j_n\}\subset\{i_0,\dots,i_n\}$ 的 $(j_0,\dots,j_n)$ 的线性组合。此外，我们有
\begin{align*}
&\partial_n(i_nf_n-\id_{K_n(I)}-h_{n-1}\partial_n)(i_0,\dots,i_n)\\
&=(i_{n-1}f_{n-1}-\id_{K_{n-1}(I)})\partial_n(i_0,\dots,i_n)-\partial_nh_{n-1}\partial_n(i_0,\dots,i_n)\\
&=(\partial_nh_{n-1}+h_{n-2}\partial_{n-1})\partial_n(i_0,\dots,i_n)-\partial_nh_{n-1}\partial_n(i_0,\dots,i_n)\\
&=0.
\end{align*}
因为 $H_n(K_\cdot(\{i_0,\dots,i_n\}))=0$，存在 $c\in K_{n+1}(\{i_0,\dots,i_n\})$ 使得
\[
\partial_{n+1}(c)=(i_nf_n-\id_{K_n(I)}-h_{n-1}\partial_n)(i_0,\dots,i_n).
\]
我们定义 $h_n(i_0,\dots,i_n)=c$。
\end{proof}

\begin{corollary}\label{it:2.3.2}
设 $X$ 是拓扑空间，$\fU=\{U_1,\dots,U_n\}$ 是 $X$ 的有限开覆盖，$\cP$ 是 $X$ 上的前层。则 $\check{H}^q(\fU,\cP)=0$ 对所有 $q\ge n$ 成立。
\end{corollary}

\begin{proof}
我们可以用交错 \v{C}ech 复形计算 \v{C}ech 上同调群。由于 $\fU$ 中只有 $n$ 个开子集，$C'^q(\fU,\cP)=0$ 对所有 $q\ge n$ 成立，断言由此得证。
\end{proof}

\begin{proposition}\label{it:2.3.3}
设 $X$ 是拓扑空间，$\fU=\{U_i\}_{i\in I}$ 是 $X$ 的开覆盖。
\begin{enumerate}[label=(\roman*),leftmargin=2.8em]
  \item $X$ 上交换群的前层范畴有足够内射对象。
  \item 交换群前层范畴上的函子 $(\check{H}^i(\fU,-))$ 构成上同调函子。
  \item 对每个内射前层 $\cI$，有 $\check{H}^i(\fU,\cI)=0$ 对每个 $i\ge 1$ 成立。
  \item $(\check{H}^i(\fU,-))$ 是泛上同调函子，$\check{H}^0(\fU,-)$ 左正合，且对任意 $i\ge 0$，$\check{H}^i(\fU,-)$ 是 $\check{H}^0(\fU,-)$ 的第 $i$ 个右导出函子。
\end{enumerate}
\end{proposition}

\begin{proof}
(i) 与 (ii) 留给读者证明。我们证明 (iii)。这将说明 $\check{H}^i(\fU,-)$ 对每个 $i\ge 1$ 都是可抹函子。于是由 \ref{it:2.1.4} 与 \ref{it:2.1.5}，(iv) 成立。

下面证明 (iii)。我们需要证 \v{C}ech 复形中的序列
\[
\prod_{i_0\in I}\cI(U_{i_0})\to\prod_{i_0,i_1\in I}\cI(U_{i_0i_1})\to\prod_{i_0,i_1,i_2\in I}\cI(U_{i_0i_1i_2})\to\cdots
\]
正合。对 $X$ 的任意开子集 $U$，在 $X$ 上定义前层 $Z_U$ 为
\[
Z_U(V)=
\begin{cases}
0, & \text{若 } V\not\subset U,\\
\bZ, & \text{若 } V\subset U.
\end{cases}
\]
则上述序列等同于
\[
\Hom\Bigl(\bigoplus_{i_0\in I}Z_{U_{i_0}},\cI\Bigr)\to
\Hom\Bigl(\bigoplus_{i_0,i_1\in I}Z_{U_{i_0i_1}},\cI\Bigr)\to
\Hom\Bigl(\bigoplus_{i_0,i_1,i_2\in I}Z_{U_{i_0i_1i_2}},\cI\Bigr)\to\cdots.
\]
由于 $\cI$ 是内射前层，只需证明序列
\[
\bigoplus_{i_0\in I}Z_{U_{i_0}}\xleftarrow{\delta_0}
\bigoplus_{i_0,i_1\in I}Z_{U_{i_0i_1}}\xleftarrow{\delta_1}
\bigoplus_{i_0,i_1,i_2\in I}Z_{U_{i_0i_1i_2}}\xleftarrow{\delta_2}\cdots
\]
是前层的正合序列，即对 $X$ 的任意开子集 $V$，序列
\[
\bigoplus_{i_0\in I}Z_{U_{i_0}}(V)\xleftarrow{\delta_0}
\bigoplus_{i_0,i_1\in I}Z_{U_{i_0i_1}}(V)\xleftarrow{\delta_1}
\bigoplus_{i_0,i_1,i_2\in I}Z_{U_{i_0i_1i_2}}(V)\xleftarrow{\delta_2}\cdots
\]
正合。下面描述
\[
\delta_n\colon\bigoplus_{i_0,\dots,i_{n+1}\in I}Z_{U_{i_0\dots i_{n+1}}}(V)\to\bigoplus_{i_0,\dots,i_n\in I}Z_{U_{i_0\dots i_n}}(V).
\]
令 $S=\{i\in I\mid V\subset U_i\}$。我们有
\[
\bigoplus_{i_0,\dots,i_n\in I}Z_{U_{i_0\dots i_n}}(V)
=\bigoplus_{V\subset U_{i_0\dots i_n}}\bZ
=\bigoplus_{(i_0,\dots,i_n)\in S^{n+1}}\bZ.
\]
于是 $\bigoplus_{i_0,\dots,i_n\in I}Z_{U_{i_0\dots i_n}}(V)$ 可与 2.3.1 的证明中定义的 $K_n(S)$ 等同。因此
\[
\delta_n\colon\bigoplus_{i_0,\dots,i_{n+1}\in I}Z_{U_{i_0\dots i_{n+1}}}(V)\to\bigoplus_{i_0,\dots,i_n\in I}Z_{U_{i_0\dots i_n}}(V)
\]
等同于 $\partial_{n+1}\colon K_{n+1}(S)\to K_n(S)$。在 2.3.1 的证明中，我们已经证明 $H^n(K_\cdot(S))=0$ 对任意 $n\ge 1$ 成立。因此序列
\[
\bigoplus_{i_0\in I}Z_{U_{i_0}}(V)\xleftarrow{\delta_0}
\bigoplus_{i_0,i_1\in I}Z_{U_{i_0i_1}}(V)\xleftarrow{\delta_1}
\bigoplus_{i_0,i_1,i_2\in I}Z_{U_{i_0i_1i_2}}(V)\xleftarrow{\delta_2}\cdots
\]
正合。
\end{proof}
\begin{tnote}
原书此处印作 $H^n(K_\cdot(S))=0$。但 $K_\cdot(S)$ 是 $\partial$ 的次数为 $-1$ 的复形，给出的是同调群 $H_n$；且 2.3.1 的证明实际证明的也是 $H_n(K_\cdot(I))=0$（$n\ge 1$）。故此处应读作 $H_n(K_\cdot(S))=0$。[依据：原书 p.170 第 3 段；原书 p.171 本段]
\end{tnote}

设 $\fU=\{U_i\}_{i\in I}$ 与 $\fV=\{V_j\}_{j\in J}$ 是 $X$ 的两个开覆盖。若存在映射 $\tau\colon I\to J$ 使 $U_i\subset V_{\tau(i)}$，则称 $\fU$ 是 $\fV$ 的\emph{加细}\idx{jiaxian}{加细}。定义
\[
\tau^*\colon C^n(\fV,\cP)\to C^n(\fU,\cP),\quad s=(s_{j_0\dots j_n})\mapsto \tau^*(s)=((\tau^*(s))_{i_0\dots i_n})
\]
为
\[
(\tau^*(s))_{i_0\dots i_n}=s_{\tau(i_0)\dots\tau(i_n)}|_{U_{i_0\dots i_n}}.
\]
则 $\tau^*\colon C^\cdot(\fV,\cP)\to C^\cdot(\fU,\cP)$ 是复形的态射，它在上同调群上诱导同态 $\tau^*\colon\check{H}^n(\fV,\cP)\to\check{H}^n(\fU,\cP)$。

\begin{proposition}\label{it:2.3.4}
记号同上。若 $\tau'\colon I\to J$ 是另一个满足 $U_i\subset V_{\tau'(i)}$ 的映射，则 $\tau^*$ 与 $\tau'^*$ 在上同调群上诱导相同的同态 $\check{H}^n(\fV,\cP)\to\check{H}^n(\fU,\cP)$。
\end{proposition}

\begin{proof}
定义
\[
h\colon C^n(\fV,\cP)\to C^{n-1}(\fU,\cP),\quad s=(s_{j_0\dots j_n})\mapsto h(s)=((h(s))_{i_0\dots i_{n-1}})
\]
为
\[
(h(s))_{i_0\dots i_{n-1}}=\sum_{k=0}^{n-1}(-1)^k s_{\tau(i_0)\dots\tau(i_k)\tau'(i_k)\dots\tau'(i_{n-1})}|_{U_{i_0\dots i_{n-1}}}.
\]
可以验证
\[
hd+dh=\tau'^*-\tau^*.
\]
所以 $\tau^*$ 与 $\tau'^*$ 同伦，从而它们在上同调群上诱导相同的同态。
\end{proof}

给定 $X$ 的两个开覆盖 $\fU$ 与 $\fV$，若 $\fU$ 是 $\fV$ 的加细，则记 $\fU\ge\fV$。对任意两个开覆盖 $\fU=\{U_i\}_{i\in I}$ 与 $\fV=\{V_j\}_{j\in J}$，开覆盖 $\fW=\{U_i\cap V_j\}_{i\in I,j\in J}$ 是 $\fU$ 与 $\fV$ 的公共加细。于是 $X$ 的开覆盖族构成有向集。由 2.3.4，只要 $\fU\ge\fV$，就有良定义的同态 $\check{H}^n(\fV,\cP)\to\check{H}^n(\fU,\cP)$。我们定义前层 $\cP$ 的第 $n$ 个\emph{\v{C}ech 上同调群} $\check{H}^n(X,\cP)$ 为
\[
\check{H}^n(X,\cP)=\dlim_{\fU}\check{H}^n(\fU,\cP).
\]

作为 2.3.3 的推论，我们有

\begin{corollary}\label{it:2.3.5}
设 $X$ 是拓扑空间。
\begin{enumerate}[label=(\roman*),leftmargin=2.8em]
  \item 交换群前层范畴上的函子 $(\check{H}^i(X,-))$ 构成上同调函子。
  \item 对每个内射前层 $\cI$，有 $\check{H}^i(X,\cI)=0$ 对每个 $i\ge 1$ 成立。
  \item $(\check{H}^i(X,-))$ 是泛上同调函子，$\check{H}^0(X,-)$ 左正合，且对每个 $i\ge 0$，$\check{H}^i(X,-)$ 是 $\check{H}^0(X,-)$ 的第 $i$ 个右导出函子。
\end{enumerate}
\end{corollary}

现在讨论层的 \v{C}ech 上同调。对 $X$ 的每个开覆盖 $\fU=\{U_i\}_{i\in I}$ 与 $X$ 上每个交换群的层 $\cF$，$\check{H}^0(\fU,\cF)$ 是同态
\[
\prod_{i\in I}\cF(U_i)\to\prod_{i,j\in I}\cF(U_i\cap U_j),\quad (s_i)\mapsto (s_j|_{U_i\cap U_j}-s_i|_{U_i\cap U_j})
\]
的核。于是有
\[
\check{H}^0(\fU,\cF)\cong \cF(X)\cong H^0(X,\cF).
\]

\begin{proposition}\label{it:2.3.6}
设 $\cF$ 是 $X$ 上的松弛层。则对 $X$ 的每个开覆盖 $\fU=\{U_i\}_{i\in I}$，有 $\check{H}^n(\fU,\cF)=0$ 对每个 $n\ge 1$ 成立。
\end{proposition}

\begin{proof}
选取层范畴中的短正合序列
\begin{equation}\label{eq:2.3.1}
0\to\cF\to\cI\to\cG\to 0
\end{equation}
使得 $\cI$ 是内射层。由 \ref{it:2.1.11}，$\cG$ 也是松弛层，且对每个开子集 $U$，序列
\[
0\to\cF(U)\to\cI(U)\to\cG(U)\to 0
\]
正合，即序列 \eqref{eq:2.3.1} 在前层范畴中正合。注意内射层 $\cI$ 也是前层范畴中的内射对象。于是 $\check{H}^n(\fU,\cI)=0$ 对每个 $n\ge 1$ 成立。与短正合序列 \eqref{eq:2.3.1} 相联系的 \v{C}ech 上同调群的长正合序列给出正合序列
\begin{center}
\begin{tikzcd}
\check{H}^0(\fU,\cI) \arrow[r] \arrow[d,equal] & \check{H}^0(\fU,\cG) \arrow[r] \arrow[d,equal] & \check{H}^1(\fU,\cF) \arrow[r] & 0\\
\cI(X) & \cG(X) & &
\end{tikzcd}
\end{center}
以及同构
\[
\check{H}^n(\fU,\cG)\cong\check{H}^{n+1}(\fU,\cF)
\]
对所有 $n\ge 1$ 成立。由于 $\cI(X)\to\cG(X)$ 是满射，上述正合序列说明 $\check{H}^1(\fU,\cF)=0$。这对每个松弛层 $\cF$ 都成立。特别地，有 $\check{H}^1(\fU,\cG)=0$。上面的同构于是说明 $\check{H}^2(\fU,\cF)=0$。对 $n$ 归纳，得到 $\check{H}^n(\fU,\cF)=0$ 对所有 $n\ge 1$ 成立。
\end{proof}

设 $\fU$ 是 $X$ 的开覆盖，$\cF$ 是 $X$ 上的层，$\cI^\cdot$ 是 $\cF$ 在层范畴中的内射分解。考虑双复形\idx{shuangfuxing}{双复形} $(K^{pq})=(C^p(\fU,\cI^q))$。我们有
\[
H^p_{\mathrm{I}}H^q_{\mathrm{II}}(K^{\cdot\cdot})\cong\check{H}^p(\fU,H^q(\cF)),
\]
其中 $H^q(\cF)$ 是由
\[
H^q(\cF)(U)=H^q(U,\cF|_U)
\]
定义的前层。

另一方面，我们有
\[
H^q_{\mathrm{I}}(K^{\cdot,p})\cong\check{H}^q(\fU,\cI^p).
\]
由于 $\cI^p$ 是内射层，由 \ref{it:2.3.6} 得 $\check{H}^q(\fU,\cI^p)=0$ 对任意 $q\ge 1$ 成立。所以
\[
H^p_{\mathrm{II}}H^q_{\mathrm{I}}(K^{\cdot\cdot})\cong
\begin{cases}
0, & \text{若 } q\ge 1,\\
H^p(\Gamma(X,\cI^\cdot))\cong H^p(X,\cF), & \text{若 } q=0.
\end{cases}
\]
于是双复形 $K^{\cdot\cdot}$ 的第二个谱序列退化，并且 $H^n(X,\cF)\cong H^n(K^{\cdot\cdot})$。因此双复形的第一个谱序列可写成
\[
E_2^{pq}=\check{H}^p(\fU,H^q(\cF))\Rightarrow H^{p+q}(X,\cF).
\]
由 \ref{it:2.2.8}(i)，该谱序列是双正则的。对 $X$ 的开覆盖 $\fU$ 取有向极限，得到双正则谱序列
\[
E_2^{pq}=\check{H}^p(X,H^q(\cF))\Rightarrow H^{p+q}(X,\cF).
\]
于是得到下述命题：

\begin{proposition}\label{it:2.3.7}
设 $\cF$ 是 $X$ 上的层，$H^q(\cF)$ 是由 $H^q(\cF)(U)=H^q(U,\cF|_U)$ 定义的前层。
\begin{enumerate}[label=(\roman*),leftmargin=2.8em]
  \item 对 $X$ 的每个开覆盖 $\fU$，有双正则谱序列
  \[
  E_2^{pq}=\check{H}^p(\fU,H^q(\cF))\Rightarrow H^{p+q}(X,\cF).
  \]
  \item 有双正则谱序列
  \[
  E_2^{pq}=\check{H}^p(X,H^q(\cF))\Rightarrow H^{p+q}(X,\cF).
  \]
\end{enumerate}
\end{proposition}

\begin{corollary}\label{it:2.3.8}
设 $\cF$ 是 $X$ 上的层，$\fU=\{U_i\}_{i\in I}$ 是 $X$ 的开覆盖。假设
\[
H^q(U_{i_0}\cap\cdots\cap U_{i_n},\cF)=0
\]
对所有 $q\ge 1$ 与所有 $i_0,\dots,i_n\in I$ 成立。则对所有的 $n$，有
\[
\check{H}^n(\fU,\cF)\cong H^n(X,\cF).
\]
\end{corollary}

\begin{proof}
我们有 $C^p(\fU,H^q(\cF))=0$ 对所有 $q\ge 1$ 成立。于是双正则谱序列
\[
E_2^{pq}=\check{H}^p(\fU,H^q(\cF))\Rightarrow H^{p+q}(X,\cF)
\]
退化。断言随即由 \ref{it:2.2.2} 得到。
\end{proof}

\begin{lemma}\label{it:2.3.9}
对 $X$ 上的每个层 $\cF$，有 $\check{H}^0(X,H^q(\cF))=0$ 对每个 $q\ge 1$ 成立。
\end{lemma}

\begin{proof}
设 $\fU=\{U_i\}_{i\in I}$ 是 $X$ 的开覆盖。注意 $\check{H}^0(\fU,H^q(\cF))$ 是同态
\[
\prod_{i\in I}H^q(\cF)(U_i)\to\prod_{i,j\in I}H^q(\cF)(U_i\cap U_j),\qquad
(s_i)\mapsto\bigl(s_j|_{U_i\cap U_j}-s_i|_{U_i\cap U_j}\bigr)
\]
的核。

设 $(s_i)$ 是 $\check{H}^0(\fU,H^q(\cF))$ 中的一个元素，其中
\[
s_i\in H^q(\cF)(U_i)=H^q(U_i,\cF|_{U_i}).
\]
取定 $\cF$ 的一个内射分解 $\cI^\cdot$。对每个 $q\ge 1$，有
\[
H^q(U_i,\cF)=\ker\bigl(\cI^q(U_i)\to\cI^{q+1}(U_i)\bigr)\big/\Image\bigl(\cI^{q-1}(U_i)\to\cI^q(U_i)\bigr).
\]
选取 $t_i\in\ker\bigl(\cI^q(U_i)\to\cI^{q+1}(U_i)\bigr)$，使其在 $H^q(U_i,\cF|_{U_i})$ 中的像为 $s_i$。由于
\[
\cI^{q-1}\to\cI^q\to\cI^{q+1}
\]
正合，存在 $U_i$ 的开覆盖 $\{U_{ij}\}_{j\in J_i}$，使得每个 $t_i|_{U_{ij}}$ 都落在 $\Image\bigl(\cI^{q-1}(U_{ij})\to\cI^q(U_{ij})\bigr)$ 中。于是 $s_i$ 在同态 $H^q(\cF)(U_i)\to H^q(\cF)(U_{ij})$ 下的像为零。令 $\fU'$ 是由 $U_{ij}$（$i\in I$，$j\in J_i$）给出的开覆盖。则 $\fU'$ 是 $\fU$ 的加细，且 $s=(s_i)$ 在同态 $\check{H}^0(\fU,H^q(\cF))\to\check{H}^0(\fU',H^q(\cF))$ 下的像为零。所以
\[
\check{H}^0(X,H^q(\cF))=\dlim_{\fU}\check{H}^0(\fU,H^q(\cF))=0.
\]
\end{proof}

\begin{corollary}\label{it:2.3.10}
对 $X$ 上的任一层 $\cF$，有 $\check{H}^0(X,\cF)\cong H^0(X,\cF)$，$\check{H}^1(X,\cF)\cong H^1(X,\cF)$ 以及 $\check{H}^2(X,\cF)\hookrightarrow H^2(X,\cF)$。
\end{corollary}

\begin{proof}
对于 \ref{it:2.3.7} (ii) 中的谱序列，当 $p<0$ 或 $q<0$ 时有
\[
E_2^{pq}=\check{H}^p(X,H^q(\cF))=0.
\]
由 \ref{it:2.2.3}，有 $E_2^{00}=\check{H}^0(X,\cF)\cong H^0(X,\cF)$，且有正合序列
\[
0\to E_2^{1,0}\to H^1\to E_2^{0,1}\to E_2^{2,0}\to H^2.
\]
由 \ref{it:2.3.9}，有 $E_2^{0,1}=\check{H}^0(X,H^1(\cF))=0$。于是 $E_2^{1,0}\cong H^1$ 且 $E_2^{2,0}\hookrightarrow H^2$，即 $\check{H}^1(X,\cF)\cong H^1(X,\cF)$ 且 $\check{H}^2(X,\cF)\hookrightarrow H^2(X,\cF)$。
\end{proof}

\begin{proposition}[Cartan 引理]\label{it:2.3.11}
设 $\fB$ 是 $X$ 的一族开子集，它们构成 $X$ 的拓扑的一个基；并设 $\cF$ 是 $X$ 上的层，使得对任意的 $U_1,\dots,U_k\in\fB$，有 $\check{H}^q(U_1\cap\cdots\cap U_k,\cF)=0$ 对所有 $q\ge 1$ 成立。则 $H^q(U_1\cap\cdots\cap U_k,\cF)=0$ 对所有 $q\ge 1$ 成立。对 $X$ 的每个开子集 $V$ 以及 $V$ 的每个由 $\fB$ 中的开子集组成的开覆盖 $\fU$，有
\[
\check{H}^n(\fU,\cF)\cong H^n(V,\cF)
\]
对所有 $n$ 成立。此外，对任意的 $n$，有
\[
\check{H}^n(V,\cF)\cong H^n(V,\cF).
\]
\end{proposition}

\begin{proof}
设 $U=U_1\cap\cdots\cap U_k$。我们对 $q\ge 1$ 归纳证明 $H^q(U,\cF)=0$。有
\[
H^1(U,\cF)\cong\check{H}^1(U,\cF)=0.
\]
设已证 $H^q(U,\cF)=0$ 对所有 $0<q<n$ 成立。则对 $U$ 的每个由 $\fB$ 中的开子集组成的开覆盖 $\fU$，有
\[
C^p(\fU,H^q(\cF))=0,\qquad
\check{H}^p(\fU,H^q(\cF))=H^p\bigl(C^\cdot(\fU,H^q(\cF))\bigr)=0,
\]
\[
\check{H}^p(U,H^q(\cF))=\dlim_{\fU}\check{H}^p(\fU,H^q(\cF))=0
\]
对所有 $0<q<n$ 成立。把 \ref{it:2.2.3} 应用于谱序列
\[
E_2^{pq}=\check{H}^p(U,H^q(\cF))\Rightarrow H^{p+q}(U,\cF),
\]
得到正合序列
\begin{center}
\begin{tikzcd}
0 \arrow[r] & E_2^{n,0} \arrow[r] \arrow[d,equal] & H^n \arrow[r] \arrow[d,equal] & E_2^{0,n} \arrow[d,equal] \\
& \check{H}^n(U,\cF) & H^n(U,\cF) & 0
\end{tikzcd}
\end{center}
其中第三个等号由 \ref{it:2.3.9} 得到。于是
\[
H^n(U,\cF)\cong\check{H}^n(U,\cF)=0.
\]
由 \ref{it:2.3.8}，对 $X$ 的每个开子集 $V$ 以及 $V$ 的每个由 $\fB$ 中的开子集组成的开覆盖 $\fU$，有 $\check{H}^n(\fU,\cF)\cong H^n(V,\cF)$ 对任意 $n$ 成立。所以
\[
\check{H}^n(V,\cF)=\dlim_{\fU}H^n(\fU,\cF)\cong H^n(V,\cF).
\]
\end{proof}

\begin{tnote}
2.3.11 证明末尾的同构式中，中间项在原书 p.176 印作 $H^n(\fU,\cF)$（有向极限号后无 check）；但由 \ref{it:2.3.8}，与 $H^n(V,\cF)$ 同构的是 \v{C}ech 上同调 $\check{H}^n(\fU,\cF)$，故此处应为 $\check{H}^n(\fU,\cF)$。[依据：原书 p.176；原书 p.174 推论 2.3.8]
\end{tnote}

设 $(X,\cO_X)$ 是环化空间。可逆 $\cO_X$-模的同构类构成一个群，其中两个元素的积由可逆层的张量积定义。平凡可逆层 $\cO_X$ 的同构类是单位元，而对任一可逆层 $\cL$，其同构类的逆是 $\mathcal{H}\mathit{om}_{\cO_X}(\cL,\cO_X)$ 的同构类。这个群称为 $X$ 的 \emph{Picard 群}\idx{picardqun}{Picard 群}，记作 $\mathrm{Pic}(X)$。

\begin{proposition}\label{it:2.3.12}
设 $(X,\cO_X)$ 是环化空间。则有典范同构
\[
H^1(X,\cO_X^{*})\cong\mathrm{Pic}(X),
\]
其中 $\cO_X^{*}$ 是由
\[
\cO_X^{*}(U)=\{s\in\cO_X(U)\mid s \text{ 是 } \cO_X(U) \text{ 中的单位}\}
\]
定义的层。它是阿贝尔群层，其群结构由环层 $\cO_X$ 的乘法诱导。
\end{proposition}

\begin{proof}
对任一开覆盖 $\fU=\{U_i\}_{i\in I}$，令 $\mathrm{Pic}(\fU)$ 是 $\mathrm{Pic}(X)$ 中由满足 $\cL|_{U_i}\cong\cO_{U_i}$ 的那些可逆层 $\cL$ 的同构类构成的子群。我们将构造一个典范同构
\[
\check{H}^1(\fU,\cO_X^{*})\cong\mathrm{Pic}(\fU).
\]
有 $\mathrm{Pic}(X)=\dlim_{\fU}\mathrm{Pic}(\fU)$。对 $X$ 的覆盖取有向极限，得到同构
\[
\check{H}^1(X,\cO_X^{*})\cong\mathrm{Pic}(X).
\]
由 \ref{it:2.3.10}，这就证明了本命题。

固定 $I$ 上的一个全序。令
\[
Z^1(\fU,\cO_X^{*})=\ker\bigl(C^1(\fU,\cO_X^{*})\to C^2(\fU,\cO_X^{*})\bigr)
\]
为
\[
\Bigl\{(s_{i_0i_1})\in\prod_{i_0<i_1}\cO_X^{*}(U_{i_0i_1})\Bigm| s_{i_1i_2}s_{i_0i_1}=s_{i_0i_2}\ \text{在 } U_{i_0i_1i_2}\ \text{上，对任意 } i_0<i_1<i_2\Bigr\},
\]
并令
\[
B^1(\fU,\cO_X^{*})=\Image\bigl(C^0(\fU,\cO_X^{*})\to C^1(\fU,\cO_X^{*})\bigr)
\]
为
\[
\Bigl\{(s_{i_0i_1})\in\prod_{i_0<i_1}\cO_X^{*}(U_{i_0i_1})\Bigm| s_{i_0i_1}=t_{i_1}t_{i_0}^{-1}\ \text{在 } U_{i_0i_1}\ \text{上，其中 } (t_{i_0})\in\prod_{i_0}\cO_X^{*}(U_{i_0})\Bigr\}.
\]
于是有
\[
\check{H}^1(\fU,\cO_X^{*})\cong Z^1(\fU,\cO_X^{*})/B^1(\fU,\cO_X^{*}).
\]
定义映射
\[
\Phi\colon Z^1(\fU,\cO_X^{*})\to\mathrm{Pic}(\fU)
\]
如下：对每个 $(s_{i_0i_1})\in Z^1(\fU,\cO_X^{*})$ 与每对 $i_0<i_1$，令
\[
\phi_{i_0i_1}\colon\cO_{U_{i_0}}|_{U_{i_0i_1}}\to\cO_{U_{i_1}}|_{U_{i_0i_1}}
\]
为把 $\cO_{U_{i_0}}|_{U_{i_0i_1}}$ 的任一截面 $s$ 映为 $\cO_{U_{i_1}}|_{U_{i_0i_1}}$ 的截面 $s s_{i_0i_1}$ 的同构。对每个 $i\in I$ 定义 $\phi_{ii}=\id_{\cO_{U_i}}$，并对每对 $i_0<i_1$ 定义 $\phi_{i_1i_0}=\phi_{i_0i_1}^{-1}$。对任意的 $i_0,i_1,i_2\in I$，在 $U_{i_0i_1i_2}$ 上有 $s_{i_1i_2}s_{i_0i_1}=s_{i_0i_2}$，所以
\[
\phi_{i_1i_2}\phi_{i_0i_1}=\phi_{i_0i_2}
\]
在 $U_{i_0i_1i_2}$ 上成立。于是可以通过 $\{\phi_{ij}\}$ 把 $\{\cO_{U_i}\}$ 粘合起来，得到一个可逆 $\cO_X$-模。定义 $\Phi((s_{i_0i_1}))\in\mathrm{Pic}(\fU)$ 为该可逆 $\cO_X$-模的同构类。可以验证
\[
\Phi\colon Z^1(\fU,\cO_X^{*})\to\mathrm{Pic}(\fU)
\]
是同态。

假设 $(s_{i_0i_1})$ 落在 $\Phi$ 的核中，并设 $\cL$ 是用 $(s_{i_0i_1})$ 给出的粘合数据按上述方式构造的可逆 $\cO_X$-模。则有一个同构 $\psi\colon\cO_X\cong\cL$。由我们的构造，对任意 $i\in I$ 有 $\cL|_{U_i}\cong\cO_{U_i}$。把这个同构与 $\psi|_{U_i}$ 复合，得到 $\cO_{U_i}$ 的一个自同构。$\cO_{U_i}$ 的任一自同构都由乘以 $\cO_X^{*}(U_i)$ 中某个元素诱导。设 $t_i\in\cO_X^{*}(U_i)$ 对应于上述自同构。可以验证，对每对 $i_0<i_1$，在 $U_{i_0i_1}$ 上有 $s_{i_0i_1}=t_{i_1}t_{i_0}^{-1}$，从而 $(s_{i_0i_1})$ 落在 $B^1(\fU,\cO_X^{*})$ 中。反之可以证明 $B^1(\fU,\cO_X^{*})$ 中每个元素都落在 $\ker\Phi$ 中。所以 $\ker\Phi=B^1(\fU,\cO_X^{*})$。

下面证明 $\Phi$ 是满射。设 $\cL$ 是可逆 $\cO_X$-模，使得对每个 $i\in I$ 有同构 $\phi_i\colon\cL|_{U_i}\cong\cO_{U_i}$。对每对 $i_1<i_2$，考虑复合
\[
\cO_{U_{i_0}}|_{U_{i_0i_1}}\xrightarrow{\;\phi_{i_0}^{-1}\;}\cL|_{U_{i_0i_1}}\xrightarrow{\;\phi_{i_1}\;}\cO_{U_{i_1}}|_{U_{i_0i_1}}.
\]
它诱导 $\cO_{U_{i_0i_1}}$ 的一个自同构。这个自同构是乘以某个 $s_{i_0i_1}\in\cO_X^{*}(U_{i_0i_1})$。于是对任意 $i_0<i_1<i_2$，在 $U_{i_0i_1i_2}$ 上有 $s_{i_0i_1}s_{i_1i_2}=s_{i_0i_2}$。所以 $(s_{i_0i_1})\in Z^1(\fU,\cO_X^{*})$。可以验证 $\Phi((s_{i_0i_1}))$ 等于 $\cL$ 的同构类。所以 $\Phi$ 是满射。

于是 $\Phi$ 诱导出同构 $\check{H}^1(\fU,\cO_X^{*})\cong Z^1(\fU,\cO_X^{*})/B^1(\fU,\cO_X^{*})$ 与 $\mathrm{Pic}(\fU)$。
\end{proof}

\begin{tnote}
2.3.12 证明中「对每对 $i_1<i_2$」在原书 p.178 如此印出，但紧随的复合式与下文诱导自同构的讨论用的都是下标 $i_0,i_1$（复合式写作 $\cO_{U_{i_0}}|_{U_{i_0i_1}}\to\cL|_{U_{i_0i_1}}\to\cO_{U_{i_1}}|_{U_{i_0i_1}}$，自同构落在 $\cO_{U_{i_0i_1}}$ 上），故此处应为「对每对 $i_0<i_1$」。[依据：原书 p.178 复合式及紧随文字]
\end{tnote}

\begin{exercises}
\item \label{ex:2.3.1}
\begin{enumerate}[label=(\arabic*),leftmargin=2.4em]
  \item 设 $\phi\colon A\to B$ 是环同态，且 $B$ 在 $A$ 上平坦。若 $a\in A$ 不是 $A$ 的零因子，则 $\phi(a)$ 不是 $B$ 中的零因子。
  \item 设 $S$ 是 $A$ 中那些不是零因子的元素构成的集合。则 $S$ 是乘性子集，且典范同态 $A\to S^{-1}A$ 是单射。设 $a\in A$。若 $\frac{a}{1}$ 在 $S^{-1}A$ 中是单位，则 $a\in S$。
  \item 设 $B$ 与 $C$ 是 $A$-代数，并设 $b\in B$。假设 $b$ 不是零因子，且 $B/(b)$ 在 $A$ 上平坦。则 $b\tensor 1$ 不是 $B\tensor_A C$ 中的零因子，且 $(B\tensor_A C)/(b\tensor 1)$ 在 $C$ 上平坦。
  \item 设 $B$ 是 $A$-代数，$b_1,b_2\in B$。假设 $b_2$ 不是零因子。则 $b_1$ 不是零因子当且仅当 $b_1b_2$ 不是零因子。假设 $b_2$ 不是零因子且 $B/(b_2)$ 在 $A$ 上平坦。则 $B/(b_1)$ 在 $A$ 上平坦当且仅当 $B/(b_1b_2)$ 在 $A$ 上平坦。
\end{enumerate}

\item \label{ex:2.3.2} 设 $X$ 是概形。对每个仿射开子集 $U=\Spec A$，设 $S_A$ 是 $A$ 中不是零因子的元素构成的集合，并令 $K(U)=S_A^{-1}A$。
\begin{enumerate}[label=(\arabic*),leftmargin=2.4em]
  \item 对 $U$ 中包含的任一仿射开子集 $V=\Spec B$，证明限制映射
  \[
  A\cong\cO_X(U)\to\cO_X(V)\cong B
  \]
  把 $S_A$ 映入 $S_B$，从而诱导同态 $K(U)\to K(V)$。于是我们得到环的前层 $U\mapsto K(U)$。设 $\cK_X$ 是与此前层相伴的层，并设 $\cK_X^{*}$ 是由可逆截面构成的子层。
  \item 证明典范态射 $\cO_X\to\cK_X$ 是单态射。

  层 $\cK_X^{*}/\cO_X^{*}$ 的整体截面称为 \emph{Cartier 除子}\idx{cartierchuzi}{Cartier 除子}。我们常把 Cartier 除子记作 $D=\{(s_i,U_i)\}$，其中 $(U_i)$ 是 $X$ 的开覆盖，$s_i\in\cK_X^{*}(U_i)$ 使得 $s_is_j^{-1}\in\cO_X^{*}(U_i\cap U_j)$。Cartier 除子的集合 $\Gamma(X,\cK_X^{*}/\cO_X^{*})$ 是阿贝尔群。若 $D'=\{(s'_j,U_j)\}$ 是另一个 Cartier 除子，则
  \[
  D+D'=\{(s_is'_j,U_i\cap U_j)\}.
  \]

  \item 对 Cartier 除子 $D=\{(s_i,U_i)\}$，定义与它相伴的层 $\cL(D)$ 为 $\cK_X$ 的如下子层：
  \[
  \cL(D)|_{U_i}=\cO_X|_{U_i}\cdot s_i^{-1}.
  \]
  证明 $\cL(D)$ 是可逆 $\cO_X$-模。设 $\cI(D)$ 是 $\cK_X$ 的如下子层：
  \[
  \cI(D)|_{U_i}=\cO_X s_i.
  \]
  则 $\cI(D)$ 是可逆 $\cO_X$-模，且 $\cI(D)\cong\cL(D)^{-1}$。

  若 Cartier 除子落在典范同态
  \[
  \Gamma(X,\cK_X^{*})\to\Gamma(X,\cK_X^{*}/\cO_X^{*})
  \]
  的像中，则称它为\emph{主除子}\idx{zhuchuzi}{主除子}。

  \item 证明 $\cL(D)$ 同构于 $\cO_X$ 当且仅当 $D$ 是主除子。

  定义 $X$ 的 \emph{Cartier 类群}\idx{cartierleiqun}{Cartier 类群} $\mathrm{CaCl}(X)$ 为上述典范同态的余核。于是得到单态射
  \[
  \mathrm{CaCl}(X)\hookrightarrow\mathrm{Pic}(X),\qquad [D]\mapsto\cL(D).
  \]

  \item 利用与短正合序列
  \[
  0\to\cO_X^{*}\to\cK_X^{*}\to\cK_X^{*}/\cO_X^{*}\to 0
  \]
  相联系的上同调群的长正合序列，证明 (4)。
  \item 设 $X$ 是整概形。证明 $\cK_X^{*}$ 是常层且是松弛层，并且有同构
  \[
  \mathrm{CaCl}(X)\cong\mathrm{Pic}(X).
  \]
\end{enumerate}

若 Cartier 除子 $D=\{(s_i,U_i)\}$ 满足 $s_i\in\cO_X(U_i)$ 对所有 $i$ 成立，则称 $D$ 为\emph{有效 Cartier 除子}\idx{youxiaocartierchuzi}{有效 Cartier 除子}。这等价于说 $\cI(D)\subset\cO_X$，即 $\cI(D)$ 是 $\cO_X$ 的理想层。这个包含诱导单态射 $\cO_X\hookrightarrow\cL(D)$。用 $s_D\in\cL(D)(X)$ 记 $1$ 的像。

\begin{enumerate}[label=(\arabic*),leftmargin=2.4em,start=7]
  \item 证明映射
  \[
  D\mapsto(\cL(D),s_D)
  \]
  定义了 $X$ 上有效 Cartier 除子的集合与下述二元组 $(\cL,s)$ 的同构类集合之间的一一对应：$\cL$ 是可逆 $\cO_X$-模，$s\in\cL(X)$ 是整体截面，且同态
  \[
  s\colon\cO_X\to\cL,\quad a\mapsto as
  \]
  是单射。

  我们称 $\Supp(\cO_X/\cI(D))$ 为 $D$ 的\emph{支集}。它也可以定义为单态射 $\cO_X\hookrightarrow\cL(D)$ 的余核的支集。

  \item 设 $g\colon X'\to X$ 是概形的平坦态射，即对每个 $x'\in X'$，$\cO_{X',x'}$ 在 $\cO_{X,g(x')}$ 上平坦。证明 $g$ 诱导层态射 $g^{*}\colon\cK_X\to g_*\cK_{X'}$。

  若 $D=\{(s_i,U_i)\}$ 是 $X$ 上的 Cartier 除子，定义 $g^{*}D$ 为 $X'$ 上的 Cartier 除子 $\{(g^{*}(s_i),g^{-1}(U_i))\}$。设 $q\colon X\to T$ 是概形态射。$X/T$ 上的\emph{相对有效 Cartier 除子}\idx{xiangduiyouxiaocartierchuzi}{相对有效 Cartier 除子}是指 $X$ 上的有效 Cartier 除子 $D=\{(s_i,U_i)\}$，使得 $X$ 的以 $\cI(D)$ 为理想层的闭子概形在 $T$ 上平坦。

  \item 证明相对有效 Cartier 除子的集合与下述二元组 $D=(\cL,s)$ 的集合之间有一一对应：$\cL$ 是可逆 $\cO_X$-模，$s\in\Gamma(X,\cL)$ 是整体截面，使得由 $s$ 诱导的同态 $s\colon\cO_X\to\cL$ 是单射且其余核在 $\cO_T$ 上平坦。
  \item 若 $D'$ 是另一个相对有效 Cartier 除子，则 $D+D'$ 也是。进一步假设 $\cI(D)\subset\cI(D')$。则 $D-D'$ 也是相对有效 Cartier 除子。
  \item 设 $T'\to T$ 是概形态射，$q'\colon X\times_T T'\to T'$ 是它的基变换。证明 $g^{*}D=\{(g^{*}(s_i),g^{-1}(U_i))\}$ 是 $X\times_T T'/T'$ 上的相对有效 Cartier 除子，其中 $g\colon X\times_T T'\to X$ 是投影。
\end{enumerate}

\item \label{ex:2.3.3} 设 $X$ 是诺特整概形，$\eta$ 是 $X$ 的一般点。假设对任意 $P\in X$（$P\ne\eta$），$P$ 是 $X$ 的闭点，且 $\cO_{X,P}$ 是离散赋值环。（例如 $X$ 是光滑代数曲线。）设 $v_P\colon K^{*}\to\bZ$ 是相应的离散赋值，其中 $K=\cO_{X,\eta}$ 是 $X$ 的函数域。用 $|X|$ 记 $X$ 的闭点集。对 $X$ 的任一开子集 $U$，定义 $U$ 上的 \emph{（Weil）除子群}\idx{weilchuziqun}{（Weil）除子群} $\mathrm{Div}(U)$ 为以 $U$ 中闭点集为基的自由阿贝尔群。闭点称为\emph{素除子}\idx{suchuzi}{素除子}。于是得到除子层 $\mathrm{Div}$：$U\mapsto\mathrm{Div}(U)$。

对任意 $f\in K^{*}$，定义与 $f$ 相伴的\emph{主除子} $(f)|_U$ 为
\[
(f)|_U=\sum_{P\in|U|}v_P(f)P.
\]
用 $K^{*}$ 记常层 $U\mapsto K^{*}$。于是有层态射 $K^{*}\to\mathrm{Div}$，由
\[
K^{*}(U)\to\mathrm{Div}(U),\quad f\mapsto (f)|_U
\]
定义。

定义\emph{类群}\idx{leiqun}{类群} $\mathrm{Cl}(X)$ 为同态 $K^{*}\to\mathrm{Div}(X)$ 的余核。
\begin{enumerate}[label=(\arabic*),leftmargin=2.4em]
  \item 证明有层的短正合序列
  \[
  0\to\cO_X^{*}\to K^{*}\to\mathrm{Div}\to 0.
  \]
  特别地，有
  \[
  \mathrm{Div}(X)\cong\Gamma(X,K^{*}/\cO_X^{*}),
  \]
  即 Weil 除子群同构于 Cartier 除子群。
  \item 证明类群之间有同构 $\mathrm{Cl}(X)\cong\mathrm{CaCl}(X)$。
\end{enumerate}
\end{exercises}


\endgroup
